% ======================================================================
% CURRÍCULO GERAL / MASTER — Gabriel Canela
% ======================================================================
% COMO USAR ESTE ARQUIVO (leia antes de editar):
%
% Este é o currículo MESTRE: contém tudo que pode ser relevante em
% alguma candidatura. Para gerar a versão de uma vaga específica:
%   1) Copie este arquivo (ex.: cv_empresa-x.tex);
%   2) Apague blocos inteiros de \entry{...}{...}{...}{...} + itemize
%      que não interessarem para a vaga (cada bloco de experiência e
%      de projeto está isolado entre comentários "% ---- INÍCIO/FIM");
%   3) Apague itens da seção "Softwares e Ferramentas" e linhas da
%      seção "Habilidades e Competências" que não forem relevantes;
%   4) Se o resultado couber em 1 página, pode reduzir a fonte para
%      10pt (troque [a4paper,11pt] por [a4paper,10pt] abaixo) e/ou
%      remover o rodapé de paginação (comente \pagestyle{fancy}).
%
% Cada bloco de Experiência/Projeto é independente — apagar um bloco
% inteiro não quebra os outros.
% ======================================================================

\documentclass[a4paper,11pt]{article}

% ---------- Codificação, idioma e tipografia ----------
\usepackage[utf8]{inputenc}
\usepackage[T1]{fontenc}          % garante extração correta de acentos (ATS/copiar-colar)
\usepackage[a4paper,left=2cm,right=2cm,top=1.6cm,bottom=1.8cm]{geometry}
\usepackage{charter}              % fonte serifada mais legível que Computer Modern, boa para impressão
\usepackage{microtype}            % melhora justificação e espaçamento fino
\usepackage{ragged2e}

% ---------- Ícones, cor e links ----------
\usepackage{fontawesome5}
\usepackage{xcolor}
\definecolor{accent}{RGB}{31,58,95}   % azul-marinho discreto — legível também em P&B
\definecolor{graytext}{RGB}{95,95,95}

\usepackage{hyperref}
\hypersetup{
    pdftitle={Gabriel Canela - Curriculo},
    pdfauthor={Gabriel Canela},
    pdfsubject={Curriculo - Engenharia de Controle e Automacao},
    colorlinks=true,
    linkcolor=accent,
    filecolor=accent,
    urlcolor=accent,
    citecolor=accent
}

% ---------- Seções ----------
\usepackage{titlesec}
\titleformat{\section}{\color{accent}\large\bfseries}{\thesection}{1em}{}[{\color{accent}\titlerule}]
\titlespacing*{\section}{0pt}{14pt}{8pt}

% ---------- Listas (espaçamento confortável, sem "parede de texto") ----------
\usepackage{enumitem}
\setlist[itemize,1]{leftmargin=1.15em, itemsep=3pt, topsep=4pt, parsep=0pt, partopsep=0pt, label=\textbullet}

% ---------- Rodapé com paginação (documento é multi-página) ----------
\usepackage{lastpage}
\usepackage{fancyhdr}
\pagestyle{fancy}
\fancyhf{}
\renewcommand{\headrulewidth}{0pt}
\renewcommand{\footrulewidth}{0.3pt}
\fancyfoot[C]{\small\color{graytext} Gabriel Canela \quad\textperiodcentered\quad Currículo \quad\textperiodcentered\quad Página \thepage\ de \pageref{LastPage}}

\setlength{\parindent}{0pt}
\setlength{\emergencystretch}{3em}

% ---------- Comando padrão para entradas (experiência/projeto/educação) ----------
% #1 = título (empresa/projeto/instituição)  #2 = data
% #3 = subtítulo (cargo/tipo)                #4 = local/observação
\newcommand{\entry}[4]{%
  \noindent\textbf{#1} \hfill \textbf{#2}\\
  \textit{#3} \hfill \textit{#4}
  \vspace{3pt}
}
\newcommand{\entrygap}{\vspace{9pt}}

\begin{document}

% ======================================================================
% CABEÇALHO
% ======================================================================
\noindent
\begin{minipage}[t]{0.52\textwidth}
\textcolor{accent}{\textbf{\Huge Gabriel Canela}}
\end{minipage}%
\begin{minipage}[t]{0.48\textwidth}
\raggedleft

\href{tel:+5519996831296}{\faPhone \space +55 (19) 99683-1296} \quad
\href{mailto:gabriel.canela.ti@gmail.com}{\faEnvelope \space gabriel.canela.ti@gmail.com}

\vspace{2pt}

\faMapMarker \space Campinas - SP, Brasil

\vspace{2pt}

\href{https://github.com/Canela-san}{\faGithub \space github.com/Canela-san} \quad
\href{https://www.linkedin.com/in/gabriel-canela-ti/}{\faLinkedin \space linkedin.com/in/gabriel-canela-ti}
\end{minipage}

\vspace{0.6em}

\begin{center}
{\Large\bfseries\color{accent} Engenharia de Controle e Automação --- Unicamp}\\[3pt]
{\normalsize\color{graytext} Sistemas Embarcados de Tempo Real \textperiodcentered\ Instrumentação \textperiodcentered\ Eletrônica e Firmware \textperiodcentered\ Automação de Processos e Dados}
\end{center}

\vspace{0.3em}

\section*{Resumo Profissional}

Estudante de Engenharia de Controle e Automação pela Unicamp, com formação técnica dupla em Informática e Informática para Internet (ETEC) e perfil focado no desenvolvimento de sistemas de precisão. Atualmente atuo como Pesquisador (Bolsista FAPESP) no desenvolvimento de instrumentação embarcada para redes elétricas, combinando modelagem de sistemas de controle, integração mecatrônica e processamento digital de sinais. Tenho experiência de ponta a ponta na condução de projetos de hardware e firmware sob demanda de clientes reais --- da concepção da PCB à entrega validada em campo --- além de vivência prática em laboratório de pesquisa (LabREI/Unicamp) na montagem e validação de conversores de potência. Proficiente em programação de baixo nível (C, C++, Assembly) e alto nível (Python, automação de scripts), com domínio de ferramentas de simulação e projeto de hardware (MATLAB/Simulink, Altium Designer, LTspice). % ATS: remova a frase abaixo se a candidatura não for para estágio/trainee
Em busca de oportunidades de estágio ou trainee em sistemas embarcados, instrumentação, automação industrial ou desenvolvimento de firmware.


\section*{Experiência}

% ---- INÍCIO: FEEC Unicamp (FAPESP) ------------------------------------
\entry{FEEC Unicamp (Projeto FAPESP)}{2025 -- Atual}{Pesquisador de Iniciação Científica}{Campinas, SP}
\begin{itemize}
\item Desenvolvimento de um sistema embarcado de medição de alta performance para identificação de supraharmônicos em redes elétricas, com taxas de amostragem na ordem de centenas de kHz.
\item Programação de firmware em C e Assembly no subsistema de tempo real (PRU-ICSS) da BeagleBone Black, controlando o barramento SPI para aquisição de dados determinística.
\item Execução de engenharia reversa e depuração física em placas de circuito impresso (PCB) customizadas de aquisição de sinal, isolando falhas de hardware sem uso de osciloscópio.
\item Reestruturação de software embarcado em ambiente Linux, com reconfiguração de Device Tree e integração de processos de persistência de dados.
\item Elaboração de setups de teste e condução de validações experimentais baseadas em processamento digital de sinais (FFT).
\end{itemize}
\entrygap
% ---- FIM: FEEC Unicamp (FAPESP) ----------------------------------------

% ---- INÍCIO: LabREI ----------------------------------------------------
\entry{LabREI (Lab. de Redes Elétricas Inteligentes) --- FEEC Unicamp}{2022 -- 2024}{Técnico de Laboratório}{Campinas, SP}
\begin{itemize}
\item Apoio técnico a pesquisadores de pós-graduação na elaboração de setups de teste e validação experimental de projetos de eletrônica de potência.
\item Integração de sistemas mecânicos e eletrônicos: cotação e escolha de componentes, usinagem/furação de painéis metálicos e execução das conexões elétricas de um conversor CA/CC bidirecional de 4 pernas de alta potência.
\item Administração da infraestrutura de TI do laboratório, gerenciando servidor local baseado em Linux Ubuntu, site institucional e atualização de certificados.
\item Treinamento de usuários na correta operação de equipamentos de medição e manutenção de hardware/software dos computadores de pesquisa.
\end{itemize}
\entrygap
% ---- FIM: LabREI --------------------------------------------------------

% ---- INÍCIO: Mecatron Jr. (NOVO) ----------------------------------------
\entry{Mecatron Projetos e Consultoria Jr.}{2022 -- 2023}{Assessor de Projetos}{Campinas, SP}
\begin{itemize}
\item Responsável de ponta a ponta pelo projeto \textbf{MecaVac} --- refrigerador inteligente para armazenamento de vacinas, desenvolvido sob demanda de um cliente externo --- da concepção ao hardware, firmware, documentação técnica e interlocução direta com o cliente.
\item Projeto completo de PCB customizada de dupla camada no Altium Designer, integrando microcontrolador ESP32, drivers de potência para atuadores e gerenciamento de energia com backup por supercapacitores.
\item Desenvolvimento do firmware embarcado em C/C++ (ESP32, PlatformIO/Arduino), implementando controle de temperatura por PWM, log de eventos em cartão microSD com RTC, monitoramento de sensores e provisionamento de rede Wi-Fi via portal cativo.
\item Condução de testes de bancada e de campo até a aprovação técnica formal do cliente, entregando um produto validado e colocado em operação.
\end{itemize}
\entrygap
% ---- FIM: Mecatron Jr. ---------------------------------------------------

\section*{Projetos e Destaques}

% ---- INÍCIO: SH-Analyzer -------------------------------------------------
\entry{SH-Analyzer: Analisador de Supraharmônicos}{2025 -- Atual}{Iniciação Científica (Bolsista FAPESP)}{github.com/Canela-san/SH-Analyzer}
\begin{itemize}
\item Desenvolvimento de um sistema híbrido de instrumentação de alto desempenho (BeagleBone Black + PCB customizada) para aquisição determinística de sinais em redes elétricas, buscando superar o limite de aprox.\ 102,4~kHz de um protótipo anterior em C puro.
\item Programação de firmware em Assembly nativo para o coprocessador de tempo real PRU-ICSS, implementando o protocolo manual do ADC ADS8688 (16~bits) e gravação em buffers \textit{ping-pong} isolados na DDR, sincronizados com o ARM por handshake em memória compartilhada.
\item Extensão da arquitetura para captura multi-canal intercalada (\textit{round-robin}) do ADC, com calibração e FFT independentes por canal nas ferramentas de pós-processamento em Python (NumPy, SciPy, Pandas, Matplotlib).
\item Depuração de hardware em bancada, isolando um bug de saturação no barramento SPI até um padrão de assimetria de tempo de subida/descida característico de um optoacoplador, por meio de um script de diagnóstico dedicado.
\end{itemize}
\entrygap
% ---- FIM: SH-Analyzer ------------------------------------------------------

% ---- INÍCIO: MecaVac (NOVO) -------------------------------------------------
\entry{MecaVac: Refrigerador Inteligente para Vacinas}{2022}{Mecatron Jr. --- Projeto sob demanda de cliente}{github.com/Canela-san/MecaVac}
\begin{itemize}
\item Projeto e roteamento de placa de circuito impresso de 2 camadas (Altium Designer) integrando o controlador ESP32-WROOM-32D, reguladores lineares, drivers de MOSFET de potência e proteção contra transientes na entrada de energia.
\item Arquitetura de firmware não bloqueante (máquina de estados, temporização por \texttt{millis()}) controlando dois módulos Peltier por PWM em múltiplos patamares de potência, a partir da leitura combinada de três sensores de temperatura (OneWire).
\item Log de eventos em cartão microSD com marcação temporal via RTC, com backup de energia baseado em supercapacitores e circuito de \textit{soft-start}, garantindo a gravação segura do evento em caso de queda de energia.
\item Provisionamento de rede Wi-Fi via portal cativo e sincronização periódica de status com um servidor remoto do cliente, mantendo o log local como redundância independente da internet.
\item Projeto validado em testes de campo e formalmente aprovado pelo cliente, do levantamento de requisitos à entrega final.
\end{itemize}
\entrygap
% ---- FIM: MecaVac ------------------------------------------------------------

% ---- INÍCIO: Custom Processor from Scratch -----------------------------------
\entry{Custom Processor from Scratch with MOSFETs}{2025 -- Atual}{Projeto de Microarquitetura e Hardware}{github.com/Canela-san/ProcessorFromScratch}
\begin{itemize}
\item Concepção de arquitetura computacional completa a partir do nível de componentes físicos, mapeando diagramas de blocos de ULA, banco de registradores e unidade de controle.
\item Modelagem e simulação de circuitos lógicos digitais no simulador \textit{Digital}, com transição planejada para validação de portas lógicas puramente baseadas em transistores MOSFET.
\item Desenvolvimento de um conjunto de instruções (ISA) e montador Assembly customizado para validação física da microarquitetura por meio da emulação de periféricos de display.
\end{itemize}
\entrygap
% ---- FIM: Custom Processor from Scratch -------------------------------------

% ---- INÍCIO: Nuvem-Privada -------------------------------------------------
\entry{Nuvem-Privada: Infraestrutura NAS Self-Hosted}{2026}{Projeto de Automação de Infraestrutura de TI}{github.com/Canela-san/Nuvem-Privada}
\begin{itemize}
\item Orquestração e containerização de um ecossistema de nuvem privada com Docker e Docker Compose, integrando Nextcloud, banco de dados MariaDB e cache Redis com verificações de saúde (\textit{health checks}) entre serviços.
\item Automação de rotinas de manutenção em segundo plano via Shell Script, com pipelines de sincronização de arquivos, limpeza de banco de dados e backups consistentes a frio.
\item Provisionamento seguro de rede via VPN mesh (Tailscale) para acesso remoto criptografado, sem exposição direta de portas do servidor à internet.
\end{itemize}
\entrygap
% ---- FIM: Nuvem-Privada -------------------------------------------------

% ---- INÍCIO: isp-benchmark-suite (NOVO) -------------------------------------------------
\entry{isp-benchmark-suite: Auditoria de Qualidade de Rede}{2026}{Projeto Pessoal}{github.com/Canela-san/isp-benchmark-suite}
\begin{itemize}
\item Suíte de scripts em Bash para testes automatizados de latência, \textit{jitter}, perda de pacotes e velocidade em rotas estratégicas (infraestrutura nacional, servidores de jogos, provedores de nuvem), com execução paralela de \texttt{ping}/\texttt{mtr} e tratamento de sinais (\textit{trap}) para limpeza segura de processos.
\item Deteção automática de degradação de link físico (ex.: queda indevida de Gigabit para Fast Ethernet) e persistência de dados estruturados em série temporal (\texttt{.jsonl}) e logs legíveis em Markdown.
\item Automação de coleta distribuída via \texttt{crontab} e módulo de análise em Python (Pandas, Seaborn, Matplotlib), gerido com \texttt{uv} (PEP 723), gerando gráficos comparativos de desempenho entre provedores.
\end{itemize}
\entrygap
% ---- FIM: isp-benchmark-suite -------------------------------------------------

% ---- INÍCIO: Motor DC PID -------------------------------------------------
\entry{Controle de Velocidade de Motor DC}{S1/2026}{Projeto Acadêmico --- Laboratório de Controle (Unicamp)}{}
\begin{itemize}
\item Projeto e sintonia de um controlador PID para controle de velocidade de um motor DC, otimizado para atingir e frear na velocidade de referência no menor tempo possível.
\item Modelagem da planta e ajuste dos ganhos do controlador pelo método do lugar das raízes (\textit{root locus}), com simulação da resposta em malha fechada em MATLAB/Simulink.
\item Análise de tempo de acomodação, sobressinal (\textit{overshoot}) e estabilidade do sistema controlado.
\end{itemize}
\entrygap
% ---- FIM: Motor DC PID -------------------------------------------------

\section*{Educação}

\entry{Unicamp --- Campinas}{2022 -- 2028 (previsto)}{Engenharia de Controle e Automação}{Campinas, SP}
\vspace{4pt}

\entry{ETEC}{2018 -- 2019}{Técnico em Informática}{}
\vspace{4pt}

\entry{ETEC}{2017 -- 2019}{Técnico em Informática para Internet}{}


\section*{Habilidades e Competências}

\begin{itemize}
\item \textbf{Controle, Automação e Modelagem:} modelagem matemática e simulação de sistemas dinâmicos; projeto, sintonia e validação de controladores em malha fechada (PID, método do lugar das raízes) em MATLAB e Simulink; simulação de circuitos digitais e transitórios em \textit{Digital} e LTspice.
\item \textbf{Sistemas Embarcados e Tempo Real:} arquiteturas de microprocessadores e microcontroladores (ARM Cortex-A8, PRU-ICSS, ESP32, RISC-V); programação de firmware em C, C++ e Assembly nativo; mapeamento de registradores, controle determinístico de \textit{jitter} e comunicação via barramentos SPI, I2C e OneWire.
\item \textbf{Eletrônica, PCB e Instrumentação:} projeto, modelagem e simulação de esquemáticos e placas de circuito impresso (PCB) multicamadas no Altium Designer, da concepção à liberação para fabricação (Gerber, BOM); drivers de potência com MOSFET/BJT, proteção contra transientes (TVS) e circuitos de backup de energia (supercapacitores, \textit{soft-start}); técnicas de isolamento galvânico (dados e potência), \textit{level shifting} e condicionamento analógico de sinais; diagnóstico físico de circuitos com osciloscópio e multímetro.
\item \textbf{Sistemas Operacionais Aplicados:} desenvolvimento embarcado em ambiente Linux (Debian/Ubuntu); configuração avançada de multiplexação de pinos (Pinmux) e Árvore de Dispositivos (Device Tree Overlay); administração de servidores locais Linux e SSH/SFTP.
\item \textbf{Programação e Ciência de Dados:} Python avançado para automação de scripts, engenharia de dados de séries temporais e pós-processamento analítico (Pandas, NumPy, SciPy, PyQt6); gerenciamento de dependências e execução de scripts isolados com \texttt{uv} (PEP 723); Processamento Digital de Sinais (DSP) focado em análise espectral via FFT e filtragem digital; consultas estruturadas em SQL.
\item \textbf{DevOps e Ferramentas de Engenharia:} versionamento e integração contínua (Git/GitHub); containerização e orquestração de microsserviços (Docker e Docker Compose) com automação em Shell Script (Bash); redes seguras com VPN mesh (Tailscale); automação e diagnóstico de rede via \texttt{cron}, \texttt{ping}, \texttt{mtr} e \texttt{jq}; ambiente de desenvolvimento VS Code.
\end{itemize}


\section*{Softwares e Ferramentas}
% Lista compacta pensada para cobertura de palavras-chave em filtros automáticos (ATS).
% Apague por categoria conforme a vaga.

\textbf{Hardware / EDA:} Altium Designer, LTspice, \textit{Digital} (simulador lógico).

\textbf{Simulação e Controle:} MATLAB, Simulink.

\textbf{Embarcados / Firmware:} PlatformIO, Arduino Framework, Texas Instruments Code Generation Tools (\texttt{clpru}), Device Tree Overlays, \texttt{remoteproc}, \texttt{config-pin}.

\textbf{Linguagens:} C, C++, Python, Assembly (PRU), SQL, \LaTeX, Bash/Shell Script.

\textbf{Python / Dados:} NumPy, Pandas, SciPy, Matplotlib, Seaborn, PyQt6, \texttt{uv}.

\textbf{DevOps / Infraestrutura:} Git, GitHub, Docker, Docker Compose, Tailscale, Linux (Ubuntu/Debian/Pop!\_OS), \texttt{cron}, \texttt{jq}, \texttt{mtr}, \texttt{speedtest-cli}, VS Code.


\section*{Cursos, Certificações e Idiomas}

\begin{itemize}
\item \textbf{Português:} nativo.
\item \textbf{Inglês --- Avançado (Nível Master Fluency):} certificação de conclusão de ciclo avançado, nível 11 --- CCAA (2015 -- 2020).
\item \textbf{Japonês --- Básico (JLPT N5):} \textit{Japanese-Language Proficiency Test} --- certificado oficial de proficiência emitido pela Japan Foundation.
\item \textbf{Formação Avançada em TI:} certificados de Qualificação Profissional em Desenvolvimento Web, Arquitetura de Redes e Banco de Dados, integrados ao Centro Paula Souza (ETEC).
\end{itemize}

\end{document}
